% =====================================================================
%  COMANDOS PERSONALIZADOS
%  --------------------------------------------------------------------
%  \complementar{...}: destaca, no corpo do texto, trechos que ainda
%  precisam ser preenchidos pelo autor (não devem permanecer na versão
%  final). Ao finalizar o TCC, basta remover os chamados deste comando.
%  =====================================================================

% Caixa de destaque para informações a serem complementadas
\newcommand{\complementar}[1]{%
  \par\medskip\noindent
  \fcolorbox{red!70!black}{red!8}{%
    \begin{minipage}{\dimexpr\linewidth-2\fboxsep-2\fboxrule\relax}
      \textbf{\textcolor{red!80!black}{[COMPLEMENTAR]:}}~#1
    \end{minipage}%
  }\par\medskip
}

% Variante para uso dentro de células de tabela, parágrafos já abertos
% ou linhas de \assinatura (sem emitir \par/medskip no início e no fim).
\newcommand{\complementarcell}[1]{%
  \fcolorbox{red!70!black}{red!8}{%
    \begin{minipage}[c]{\dimexpr\linewidth-2\fboxsep-2\fboxrule\relax}
      \small\textbf{\textcolor{red!80!black}{[COMPLEMENTAR]:}}~#1
    \end{minipage}%
  }%
}

% Etiqueta compacta para células de tabela muito estreitas (sem caixa).
\newcommand{\complementartag}{%
  \textcolor{red!80!black}{\footnotesize\textbf{[COMPLEMENTAR]}}%
}

% Nota editorial genérica (subtítulo em versão de rascunho)
\newcommand{\rascunho}[1]{%
  \par\medskip\noindent
  \fcolorbox{gray}{gray!8}{%
    \begin{minipage}{\dimexpr\linewidth-2\fboxsep-2\fboxrule\relax}
      \textbf{\textcolor{gray}{[RASCUNHO]:}}~#1
    \end{minipage}%
  }\par\medskip
}

% Atalho para código com fonte monoespaçada
\newcommand{\code}[1]{\texttt{#1}}

% Nome realçado do projeto
\newcommand{\argus}{\textbf{Argus Cyclist}}
